% !TEX program = pdflatex
\documentclass[11pt]{article}
\usepackage[a4paper,margin=1in]{geometry}
\usepackage{amsmath,amssymb,mathtools,amsfonts}
\usepackage{graphicx}
\usepackage{hyperref}
\usepackage[nameinlink]{cleveref}
\usepackage{booktabs}
\usepackage{microtype}
\usepackage{enumitem}
\usepackage{listings}
\usepackage{xcolor}
\usepackage{algorithm}
\usepackage{algpseudocode}
\lstset{basicstyle=\ttfamily\small,breaklines=true,columns=fullflexible,frame=single,showstringspaces=false}
\hypersetup{colorlinks=true,linkcolor=blue,citecolor=blue,urlcolor=blue}
\title{A Stable PDE-RNN with Sparse Modular Memory: Specification, Analysis, and Implementation Notes}
\author{Ryan O. Van Gelder \\ \\ Citizen Gardens}
\date{October 23, 2025}
\begin{document}
\maketitle
\begin{abstract}
We formalize a compact recurrent architecture that couples a provably contractive neural update with a sparse modular memory (SMM) and a simple frequency-domain edge operator. We provide sufficient stability conditions, complexity bounds for memory retrieval, and low-rank parameterizations that cut parameters while preserving stability. We include minimal PyTorch-style code snippets and an evaluation protocol. Claims are limited to empirical performance on standard tasks.
\end{abstract}

\section{Notation}
Time index $t=1,\dots,T$; step $\Delta t>0$. Inputs $x_t\in\mathbb R^{d}$; hidden $h_t\in\mathbb R^{n}$; outputs $y_t\in\mathbb R^{o}$. Parameters: $A\in\mathbb R^{n\times n}$, $B\in\mathbb R^{n\times n}$, $C\in\mathbb R^{n\times n}$, $U\in\mathbb R^{n\times d}$, $W\in\mathbb R^{o\times n}$. Nonlinearity $\sigma$ is 1-Lipschitz. $\|\cdot\|$ is spectral norm, $\rho(\cdot)$ spectral radius.

\section{Model}
\paragraph{Continuous-time dynamics}
\begin{equation}
\frac{\mathrm d h}{\mathrm d t} = A h + B\,\sigma(C h + U x) - \gamma h, \qquad \gamma\ge 0.
\end{equation}
\paragraph{Discrete cell (Euler)}
\begin{equation}\label{eq:r1}
\boxed{\ h_{t+1}=h_t + \Delta t\big[(A-\gamma I) h_t + B\,\sigma(C h_t + U x_t)\big]\ }
\end{equation}
\begin{equation}\label{eq:r2}
\boxed{\ y_t = W h_t.\ }
\end{equation}

\section{Stability}
Define $\Phi(h)=h+\Delta t[(A-\gamma I)h+B\,\sigma(C h+Ux)]$. If $\sigma$ is 1-Lipschitz then
\begin{equation}\label{eq:contract}
L \;\triangleq\; \operatorname{Lip}(\Phi) \le \big\|I+\Delta t(A-\gamma I)\big\| + \Delta t\,\|B\|\,\|C\|.
\end{equation}
A sufficient condition for contraction is $L<1$. A simple recipe sets $A=-\alpha I$ and $\gamma=0$, giving
\begin{equation}
L = \big|1-\Delta t\,\alpha\big| + \Delta t\, b c,\qquad b=\|B\|,\ c=\|C\|.
\end{equation}
Choose $\Delta t\,\alpha\le 1$. If $\alpha>bc$, then \eqref{eq:contract} yields a contraction. The near-optimal step under this bound is $\Delta t^* = 1/\alpha$ with $L=bc/\alpha$. Target $L\in[0.83,0.95]$ via $\alpha=(1+\varepsilon)bc$, $\varepsilon\in[0.05,0.2]$. Backpropagated gradients satisfy $\|\partial h_{t+\tau}/\partial h_t\|\le L^{\tau}$.

\paragraph{Non-diagonal extension} Let $A=-\alpha I + S$ with $S=-S^\top$ and $\|S\|\le s$. Then a sufficient bound is
\begin{equation}
L \le \big|1-\Delta t\,\alpha\big| + \Delta t\,(bc + s) < 1.
\end{equation}
We recommend $s\le 0.1\,\alpha$.

\paragraph{Implicit step} Implicit Euler enlarges the stability region:
\begin{equation}
 h_{t+1}=(I-\Delta t(A-\gamma I))^{-1}\big[h_t+\Delta t\,B\,\sigma(C h_t + U x_t)\big],
\end{equation}
solved approximately by CG (5--10 iterations, warm start from $h_t$).

\section{Parameterization and Efficiency}
Keep $A=-\alpha I+S$ with low-rank skew $S=Q R^\top - R Q^\top$ and clamp $\|S\|\le s$. Factor $B=L_B R_B^\top$, $C=L_C R_C^\top$ with ranks $r_B,r_C$. Parameter count drops from $2n^2$ to $2n(r_B+r_C)$. Example: $n=512$, $r_B=r_C=64$ yields $2\cdot 512\cdot 128=131{,}072$ params vs $2\cdot 512^2=524{,}288$ (\(\approx\)75\% reduction).

\section{Sparse Modular Memory (SMM)}
Choose dictionary $M\in\mathbb R^{D\times K}$ with unit-norm columns $m_k$. A concept is an index set $S\subseteq[K]$, $|S|=r$. Encode by superposition: $z=\sum_{k\in S} m_k = M s$ with $s\in\{0,1\}^K$, $\|s\|_0=r$.

\paragraph{Threshold retrieval with coherence} Let $u=M^\top z$ and coherence $\mu=\max_{i\ne j}|m_i^\top m_j|$. With noise $\eta$ in $z$,
\begin{align}
\min_{k\in S} u_k &\ge 1-(r-1)\mu - \|\eta\|,\\
\max_{\ell\notin S} u_\ell &\le r\mu + \|\eta\|.
\end{align}
Hence a valid threshold interval exists if $(2r-1)\mu + 2\|\eta\|<1$; pick $\tau$ in that interval and accept indices with $u_k\ge\tau$ \cite{Tropp2004GreedIsGood}.

\paragraph{Random dictionary sizing} For i.i.d. subgaussian columns, $\mu \lesssim c\sqrt{(\log K)/D}$ with high probability; a practical rule is $D \ge C r^2 \log K$ with $C\in[8,16]$ \cite{JL1984,JL_Tutorial2021}.

\paragraph{Complexity and accelerations} Na"{i}ve read is $O(DK)$. Reduce cost with: (i) HNSW ANN prefilter, expected $O(\log K + BD)$ for candidate set size $B$ \cite{Malkov2016HNSW}; (ii) sketch-then-select using CountSketch/JL projections with high-recall guarantees \cite{Charikar2002CountSketch,JL1984}.

\section{Memory--RNN Interface}
Write: $m_{t+1}=\lambda m_t + z_t$, $\lambda\in[0,1)$. Read: $r_t=M^\top m_t$. Gated injection keeps the stability bound unchanged if $r_t$ is treated constant at step $t$:
\begin{align}
 g_t &= \sigma_g(G h_t + H x_t),\\
 \tilde u_t &= C h_t + U x_t + P r_t,\\
 h_{t+1} &= h_t + \Delta t\big[(A-\gamma I)h_t + g_t\odot B\,\sigma(\tilde u_t)\big].
\end{align}

\section{Edge Operator (FFT)}
For image $I\in\mathbb R^{H\times W}$: $F=\operatorname{FFT2}(I)$, apply high-pass mask $\mathcal M$ outside radius $\rho$ around DC, then $E=\operatorname{IFFT2}(\mathcal M\odot F)$ \cite{CooleyTukey1965}. Complexity $O(HW\log(HW))$.

\section{Training Objectives}
Sequence loss $\mathcal L_{\text{task}} = \tfrac1T\sum_t \ell(Wh_t, y_t^{\text{true}})$.
Optional memory discrimination:
\begin{equation}
\mathcal L_{\text{mem}} = \frac{1}{|S_t|}\sum_{k\in S_t} \max(0,\tau-u_k) + \frac{1}{|N_t|}\sum_{k\in N_t} \max(0, u_k-\tau).
\end{equation}
Total loss $\mathcal L=\mathcal L_{\text{task}}+\lambda_{\text{mem}}\,\mathcal L_{\text{mem}}$.

\section{Implementation Snippets}
\subsection*{Recurrent step (PyTorch-like)}
\begin{lstlisting}[language=Python]
# h: (n,), x: (d,)
# Parameters: A, B, C, U, W; scalars dt, gamma
h = h + dt * (A @ h + B @ sigma(C @ h + U @ x) - gamma * h)
y = W @ h
\end{lstlisting}

\subsection*{Two-vector spectral clamp (power iteration)}
\begin{lstlisting}[language=Python]
@torch.no_grad()
def clamp_spectral(W, target=1.0, iters=5):
    # W: (m, n)
    u = F.normalize(torch.randn(W.size(0), 1, device=W.device), dim=0)
    v = torch.randn(W.size(1), 1, device=W.device)
    for _ in range(iters):
        v = F.normalize(W.t() @ u, dim=0)
        u = F.normalize(W @ v, dim=0)
    sigma = (u.t() @ (W @ v)).abs().item()
    scale = target / max(sigma, 1e-6)
    W.mul_(scale)
\end{lstlisting}

\subsection*{SMM write/read}
\begin{lstlisting}[language=Python]
# Offline: M in R^{D x K} with unit-norm columns
m = torch.zeros(D)

def write(indices, M, m, lambda_):
    z = M[:, indices].sum(dim=1)
    return lambda_ * m + z

def read(M, m, tau=0.5):
    u = M.t() @ m
    return (u >= tau), u
\end{lstlisting}

\subsection*{ANN prefilter skeleton}
\begin{lstlisting}[language=Python]
# Use FAISS or HNSW for large K, then exact re-score candidates
# candidates = index.search(m[None, :].cpu().numpy(), B)[1][0]
# u_c = (M[:, candidates].t() @ m)
\end{lstlisting}

\section{Evaluation Protocol}
\textbf{E1: Modular memory toy.} r-hot concepts over K with noise. Metrics: top-k recall of $S$, FPR vs $(\mu,r,D)$. Pass: recall $\ge 95\%$, FPR $\le 5\%$.

\textbf{E2: PDE-RNN vs GRU/LSTM.} Copying, addition, permuted-MNIST, long-range retrieval. Match parameters and budget. Metrics: accuracy, CE/bpc, gradient norms, wall-clock. Pass: win/tie on $\ge2$ tasks at $\le$ compute.

\textbf{E3: FFT edges.} Natural images + synthetic noise. Baselines: Sobel, Canny, high-pass FIR. Metrics: PSNR/SSIM vs reference, runtime. Pass: match FIR quality and runtime class.

\section{Training Schedule and Diagnostics}
Target contraction $L=0.95$ initially. Set $\alpha=(1+\varepsilon)bc$ with $\varepsilon=0.05$, take $\Delta t=1/\alpha$. Warm up the spectral clamps on $B,C$ from 0.8 to 1.0 over first 10\% of steps. Global grad-norm clip 1.0. Monitor stability margin $\delta=1-L$, retrieval gap $\Delta=\min_{k\in S} u_k-\max_{\ell\notin S} u_\ell$, and memory utilization.

\section{Related Work}
Connections to continuous-depth models and stability-focused RNNs include Neural ODEs \cite{Chen2018NeuralODE} and antisymmetric RNNs \cite{Chang2019ASRNN}. Spectral normalization provides efficient operator-norm control \cite{Miyato2018SN}. Coherence-based sparse recovery guarantees underlie our retrieval bound \cite{Tropp2004GreedIsGood}. Dimensionality reduction and sketching tools include the Johnson--Lindenstrauss lemma \cite{JL1984,JL_Tutorial2021} and CountSketch \cite{Charikar2002CountSketch}. HNSW offers fast high-recall ANN \cite{Malkov2016HNSW}. FFT edge extraction follows the classic Cooley--Tukey algorithm \cite{CooleyTukey1965}.

\section{Reproducibility Checklist}
Fixed RNG seeds; log $\|A\|,\|B\|,\|C\|$ after clamp; save configs and environment hashes; report variance over 5 runs.

\begin{thebibliography}{99}
\bibitem{Chen2018NeuralODE} R.~T.~Q. Chen, Y. Rubanova, J. Bettencourt, and D.~K. Duvenaud, ``Neural Ordinary Differential Equations,'' in \emph{NeurIPS}, 2018. \url{https://arxiv.org/abs/1806.07366}.

\bibitem{Chang2019ASRNN} B. Chang, M. Chen, E. Haber, and E.~H. Chi, ``AntisymmetricRNN: A Dynamical System View on Recurrent Neural Networks,'' 2019. \url{https://openreview.net/pdf?id=ryxepo0cFX}.

\bibitem{Miyato2018SN} T. Miyato, T. Kataoka, M. Koyama, and Y. Yoshida, ``Spectral Normalization for Generative Adversarial Networks,'' 2018. \url{https://arxiv.org/abs/1802.05957}.

\bibitem{Tropp2004GreedIsGood} J.~A. Tropp, ``Greed is Good: Algorithmic Results for Sparse Approximation,'' 2004. Preprint. \url{https://tropp.caltech.edu/papers/Tro04-Greed-Good-preprint.pdf}.

\bibitem{JL1984} W.~B. Johnson and J. Lindenstrauss, ``Extensions of Lipschitz mappings into a Hilbert space,'' 1984.

\bibitem{JL_Tutorial2021} B. Ghojogh et al., ``Johnson--Lindenstrauss Lemma, Linear and Nonlinear Random Projections: A Tutorial,'' 2021. \url{https://arxiv.org/pdf/2108.04172}.

\bibitem{Charikar2002CountSketch} M. Charikar, K. Chen, and M. Farach-Colton, ``Finding Frequent Items in Data Streams,'' 2002. \url{https://dl.acm.org/doi/10.5555/646255.684566}.

\bibitem{Malkov2016HNSW} Y.~A. Malkov and D.~A. Yashunin, ``Efficient and Robust Approximate Nearest Neighbor Search Using Hierarchical Navigable Small World Graphs,'' 2016. \url{https://arxiv.org/abs/1603.09320}.

\bibitem{CooleyTukey1965} J.~W. Cooley and J.~W. Tukey, ``An Algorithm for the Machine Calculation of Complex Fourier Series,'' \emph{Mathematics of Computation}, 1965.
\end{thebibliography}

\end{document}
